\section{Extremal discs and boundary regularity}
\label{sec:extremal-discs}

Our goal is to maximize the disc functional and obtain enough boundary
regularity to vary the whole disc. Hyperbolicity first gives compactness
on interior subdiscs; an area penalty will control the boundary tail.
Throughout this section, $k$ is the fixed K\"ahler metric on $X$.
We write $|f'(z)|_k^2=k(f'(z),f'(z))$ and use the path distance of
the associated Riemannian metric when discussing uniform convergence.
The fixed numerical factor relating real and complex tangent norms will
be absorbed into constants.

\begin{lemma}[Derivative bound and normality]
\label{ex:brody}
There is a constant $D>0$, depending only on $(X,k)$, such that every
holomorphic map $f\colon\D\to X$ satisfies
\[
 d(f)=|f'(0)|_k^2\leq D.
\]
Consequently,
\begin{equation}
 |f'(z)|_k\leq\frac{\sqrt D}{1-|z|},\qquad z\in\D.
 \label{eq:brody-derivative}
\end{equation}
Every sequence of such maps has a subsequence converging uniformly on
compact subsets of $\D$ to a holomorphic map into $X$. In local target
coordinates, all derivatives converge uniformly on smaller compact
subsets as well.
\end{lemma}

\begin{proof}
We give Brody's rescaling argument in the normalization needed here
\cite{Brody1978}.
Suppose first that there are maps $f_j\colon\D\to X$ with
$|f_j'(0)|_k\to\infty$. Choose $z_j$ at which the continuous function
\[
 z\longmapsto (1/2-|z|)|f_j'(z)|_k
\]
attains its maximum on $\{|z|\leq1/2\}$. For all sufficiently large $j$
this maximum is positive, so $|z_j|<1/2$. Put
\[
 \lambda_j=|f_j'(z_j)|_k,
 \qquad R_j=(1/2-|z_j|)\lambda_j,
 \qquad g_j(w)=f_j(z_j+w/\lambda_j).
\]
The map $g_j$ is defined on $\{|w|<R_j\}$, and
\[
 R_j\geq\tfrac12|f_j'(0)|_k\longrightarrow\infty,
 \qquad |g_j'(0)|_k=1.
\]
If $|w|<R_j/2$, then
\[
 1/2-|z_j+w/\lambda_j|
 \geq\tfrac12(1/2-|z_j|).
\]
The maximizing property of $z_j$ therefore gives
$|g_j'(w)|_k\leq2$ on this smaller disc.

For every fixed radius these maps are consequently equicontinuous and
take values in the compact metric space $X$. The Arzel\`a--Ascoli theorem
and a diagonal subsequence give a locally uniform limit
$g\colon\C\to X$. To verify holomorphy, choose a point of the source and
a target coordinate chart containing its image under $g$. On a
sufficiently small source neighborhood, the limit and all sufficiently
late maps take values in that chart. Their coordinate representations
converge uniformly on smaller neighborhoods. The usual theorem on
uniform limits of holomorphic functions, followed by the Cauchy
integral formula for derivatives, shows that $g$ is holomorphic and
$|g'(0)|_k=1$. This contradicts the hypothesis on entire maps. The
central derivative bound follows.

For $z\in\D$, apply this bound to
$w\mapsto f(z+(1-|z|)w)$. Its central derivative is
$(1-|z|)f'(z)$, proving Equation~\eqref{eq:brody-derivative}. The estimate gives
equicontinuity on every compact source disc. A second application of
Arzel\`a--Ascoli and the same coordinate argument proves normality and
convergence of derivatives. Increasing $D$ if necessary makes it
strictly positive.
\end{proof}

We next fix a point $x\in X$, a smooth K\"ahler form $h$, a smooth
closed real $(1,1)$-form $q$, and a number $\delta>0$. All constants in
the next two results may depend on these choices. Finite area can be
measured using either $h$ or $k$, since the two metrics are uniformly
comparable on $X$. We consider
\begin{equation}
 J_{h,q,\delta}(f)=d(f)+P_q(f)-\delta A_h(f)
 \label{eq:disc-functional}
\end{equation}
on the finite-area holomorphic discs with $f(0)=x$.

\begin{proposition}[Existence of a maximizing disc]
\label{ex:maximizer}
The functional in Equation~\eqref{eq:disc-functional} has finite
supremum, attained by a finite-area holomorphic disc centered at $x$.
\end{proposition}

\begin{proof}
Choose $C_q\geq0$ such that
\[
 |q(v,v)|\leq C_q h(v,v)
\]
for every tangent vector $v$. Select $r_0\in[1/2,1)$ so close to $1$
that $2C_q\log(1/r_0)\leq\delta/2$. For a disc $f$ write
\[
 a_f(z)=h(f'(z),f'(z)),\qquad b_f(z)=q(f'(z),f'(z)).
\]
On the outer annulus $r_0<|z|<1$, the function
\[
 T_f(z)=\delta a_f(z)-2\log(1/|z|)b_f(z)
\]
satisfies
\begin{equation}
 T_f(z)\geq\frac\delta2 a_f(z)\geq0.
 \label{eq:negative-tail-density}
\end{equation}
On the inner disc, Lemma~\ref{ex:brody} and comparison of $h$ with $k$
give a uniform bound for $a_f$ and $|b_f|$. Since
$\log(1/|z|)$ is integrable there, there is a constant $K$ independent
of $f$ such that
\[
 \left|\frac1\pi\int_{|z|<r_0}
 \bigl(2\log(1/|z|)b_f(z)-\delta a_f(z)\bigr)\,dS\right|
 \leq K.
\]
Consequently,
\begin{equation}
 J_{h,q,\delta}(f)
 \leq D+K-\frac\delta{2\pi}\int_{r_0<|z|<1}a_f\,dS.
 \label{eq:maximizing-area-bound}
\end{equation}
In particular, the supremum is finite. It is nonnegative because the
constant disc has value zero.

Choose a maximizing sequence $(f_j)$ with $J_{h,q,\delta}(f_j)\geq-1$.
Equation~\eqref{eq:maximizing-area-bound} bounds its outer-annulus
areas uniformly, and the compact-interior derivative estimate bounds
its inner areas. Thus
\[
 \sup_j A_h(f_j)<\infty.
\]
By Lemma~\ref{ex:brody}, after passing to a subsequence the maps converge
locally uniformly to a holomorphic disc $f$ centered at $x$. Derivative
convergence on compact subsets gives pointwise convergence of $a_{f_j}$
and $b_{f_j}$. Fatou's lemma gives
$A_h(f)\leq\liminf_j A_h(f_j)<\infty$.

We verify the required upper semicontinuity of $J$. Its central term
converges by derivative convergence. Its inner-disc integral converges
by dominated convergence, using the uniform bound above together with
the integrable logarithmic weight. On the outer annulus, the
nonnegative functions $T_{f_j}$ converge pointwise to $T_f$, so Fatou's
lemma yields
\[
 -\frac1\pi\int_{r_0<|z|<1}T_f\,dS
 \geq\limsup_j
 \left(-\frac1\pi\int_{r_0<|z|<1}T_{f_j}\,dS\right).
\]
All these integrals are finite. Indeed, finite area controls the
unweighted integrals and the logarithmic weight is bounded on the
outer annulus. Combining the three terms proves
\[
 J_{h,q,\delta}(f)\geq\limsup_jJ_{h,q,\delta}(f_j),
\]
which proves attainment.
\end{proof}

Recall that $A_{q,s}(f)$ is the area integral restricted to $|z|<s$.
The form $q$ need not be positive, so this quantity can have either
sign. Nevertheless, $|q|\leq C_qh$ implies that $A_{q,s}(f)$ is bounded
in $s$ and tends to $A_q(f)$ as $s$ tends to $1$.
Maximality now improves the finite area of the disc to a quantitative
tail estimate. This is the additional control needed at the boundary.

\begin{lemma}[Dilation and area tails]
\label{ex:dilation}
If $f$ is a maximizer in Proposition~\ref{ex:maximizer}, then, for
$0<r<1$,
\begin{equation}
 \begin{split}
 0\leq\delta\bigl(A_h(f)-A_{h,r}(f)\bigr)
 &\leq(1-r^2)d(f)
       +2\int_r^1 A_{q,s}(f)\,\frac{ds}{s}.
 \end{split}
 \label{eq:tail}
\end{equation}
In particular,
\begin{equation}
 A_q(f)\geq-d(f),\qquad
 A_k(f)-A_{k,r}(f)\leq C_f(1-r)\quad(1/2\leq r<1)
 \label{eq:extremal-area-tail}
\end{equation}
for some finite constant $C_f$.
\end{lemma}

\begin{proof}
Let $f_r(z)=f(rz)$. A change of variables gives
\[
 d(f_r)=r^2d(f),\qquad
 A_h(f_r)=A_{h,r}(f),\qquad
 A_{q,s}(f_r)=A_{q,rs}(f).
\]
The identity $\log(1/|z|)=\int_{|z|}^1ds/s$ and Fubini's theorem give
\[
 P_q(f)=2\int_0^1 A_{q,s}(f)\,\frac{ds}{s}.
\]
Fubini is valid also for this signed integrand: near the origin the
coordinate derivatives of $f$ are bounded, while away from the
origin the logarithmic weight is bounded and $|q(f',f')|$ is controlled
by the finite $h$-area. Applying the same identity to $f_r$ gives
\begin{equation}
 P_q(f)-P_q(f_r)
   =2\int_r^1 A_{q,s}(f)\,\frac{ds}{s}.
 \label{eq:dilation-poisson}
\end{equation}
Since $f_r$ has the same center and finite area, maximality gives
$J_{h,q,\delta}(f)\geq J_{h,q,\delta}(f_r)$. Substituting the three
identities proves Equation~\eqref{eq:tail}; its leftmost inequality uses the
positivity of $h$.

Divide the rightmost expression in Equation~\eqref{eq:tail} by $1-r$. As
$r\uparrow1$, it tends to $2d(f)+2A_q(f)$. It is nonnegative for every
$r$, and hence $A_q(f)\geq-d(f)$. For $r\geq1/2$, the original
right-hand side of Equation~\eqref{eq:tail} is at most
\[
 2d(f)(1-r)+4C_qA_h(f)(1-r).
\]
This bounds the $h$-area tail by a constant times $1-r$. Uniform
comparison of $k$ and $h$ then gives the second assertion in
Equation~\eqref{eq:extremal-area-tail}.
\end{proof}

The tail estimate supplies the last step of this section: continuity
on the closed disc and a Sobolev exponent strictly above two. These
are precisely the hypotheses used to construct the frames and
variations in Section~\ref{var:section}.
\begin{proposition}[Boundary regularity of maximizing discs]
\label{ex:regularity}
Every disc maximizing the functional in Equation~\eqref{eq:disc-functional}
extends continuously
to $\overline\D$ and belongs to $W^{1,p}(\D,X)$ for every $2<p<4$.
In particular, it has continuous $W^{1,3}$ regularity up to the
boundary. More precisely, there is a constant $C_f$ such that
\begin{equation}
 |f'(z)|_k\leq C_f(1-|z|)^{-1/2}
 \label{eq:boundary-derivative}
\end{equation}
for $z$ sufficiently close to $\partial\D$.
\end{proposition}

\begin{proof}
We first make precise the uniform choice of target charts for small
source discs. Choose finitely many nested coordinate neighborhoods
\[
 V_\alpha\Subset W_\alpha\Subset U_\alpha
\]
such that the $V_\alpha$ cover $X$ and each $U_\alpha$ is a holomorphic
coordinate chart. By compactness, there is a number $\rho>0$ such that
if $y\in V_\alpha$, its metric ball of radius $\rho$ is contained in
$W_\alpha$. On each $\overline{W_\alpha}$, the metric $k$ and the
Euclidean coordinate metric are uniformly comparable; one comparison
constant suffices for this finite collection.

For $z\in\D$ put $t=1-|z|$. Let $0<c<1/4$, to be fixed independently
of $z$ and $f$. For $w\in B(z,ct)$, the straight segment from $z$ to
$w$ stays at source distance at least $(1-c)t$ from $\partial\D$.
Equation~\eqref{eq:brody-derivative} therefore bounds the length of
its image by
\[
 C\sqrt D\,\frac{c}{1-c},
\]
where $C$ accounts only for the convention relating real and complex
tangent norms. Choose $c$ small enough that this bound is less than
$\rho$. Selecting $\alpha$ with $f(z)\in V_\alpha$, we obtain
\[
 f\bigl(B(z,ct)\bigr)\subset W_\alpha.
\]

In these coordinates the derivative of $f$ is a vector of holomorphic
functions. Its squared Euclidean norm is subharmonic. The mean-value
inequality and the uniform metric comparisons give
\[
 |f'(z)|_k^2
 \leq\frac{C_1}{t^2}
       \int_{B(z,ct)}|f'(w)|_k^2\,dS(w).
\]
Every point of $B(z,ct)$ has modulus greater than
$1-(1+c)t$. For $t$ sufficiently small, the area-tail estimate in
Equation~\eqref{eq:extremal-area-tail} consequently gives
\[
 \int_{B(z,ct)}|f'(w)|_k^2\,dS(w)
 \leq\pi\bigl(A_k(f)-A_{k,1-(1+c)t}(f)\bigr)
 \leq C_2t.
\]
This proves Equation~\eqref{eq:boundary-derivative}.

For $r<s<1$ sufficiently close to $1$ and $\zeta\in\partial\D$,
integration along the radial segment gives
\[
 \operatorname{dist}_k\bigl(f(r\zeta),f(s\zeta)\bigr)
 \leq C_3\int_r^s(1-u)^{-1/2}\,du
 \leq2C_3\sqrt{1-r}.
\]
The estimate is uniform in $\zeta$. Since $X$ is complete, each radial
path has a limit, and the continuous maps
$\zeta\mapsto f(r\zeta)$ converge uniformly as $r\uparrow1$. Their
limit is therefore continuous on $\partial\D$. The same uniform
estimate, together with continuity of this boundary map, proves that
the resulting extension of $f$ is continuous on all of $\overline\D$.

It remains to improve the integrability of the derivative. For all
sufficiently large integers $j$, let
\[
 E_j=\{z\in\D:2^{-j-1}<1-|z|<2^{-j}\}.
\]
The area-tail estimate and Equation~\eqref{eq:boundary-derivative} imply
\[
 \int_{E_j}|f'|_k^2\,dS\leq C_4\,2^{-j},
 \qquad
 \sup_{E_j}|f'|_k\leq C_5\,2^{j/2}.
\]
Thus, for $2<p<4$,
\[
 \begin{split}
 \int_{E_j}|f'|_k^p\,dS
 &\leq\bigl(\sup_{E_j}|f'|_k\bigr)^{p-2}
              \int_{E_j}|f'|_k^2\,dS\\
 &\leq C_p\,2^{-j(4-p)/2}.
 \end{split}
\]
This is summable in $j$. On the remaining compact interior disc the
derivative is bounded, so $|f'|_k\in L^p(\D)$.

For completeness, the asserted Sobolev regularity is understood in
local target coordinates. Continuity on $\overline\D$ allows a finite
cover of the closed source disc by small neighborhoods whose images
lie in relatively compact target charts. In each such chart, the
coordinate functions are bounded, are smooth in the source interior,
and have first derivatives in $L^p$ by the estimate just proved and
metric comparison. These classical derivatives are their weak
derivatives, as is seen by testing against compactly supported smooth
functions in the source interior. Hence the coordinate functions lie
in $W^{1,p}$ up to the boundary, with the continuous boundary values
already constructed. This proves the proposition.
\end{proof}

\begin{remark}
The constants in Proposition~\ref{ex:regularity} need not be uniform
in $h$, $q$, or $\delta$. What will be used is the regularity of each
individual maximizing disc, to construct variations with unrestricted
boundary values. The constant $D$ in Lemma~\ref{ex:brody}, in contrast,
depends only on the fixed metric $k$ and the compact manifold $X$.
\end{remark}
